\section{Introduction}\label{sec:introduction}
Fix an odd integer \(L\ge5\), a power \(M=L^m\ge L^2\), and
\(0<\alpha<1\). At level \(k\ge0\), set
\(\eta=L^{-k}\) and \(N=\eta^{-1}\).
We use real periodic cochains on
\(T_{\eta,\ell}=\eta\Z^3/\ell\Z^3\), with pairing
\(\eta^3\sum\) on each component. Differences are normalized by
the mesh:
\[
 \partial_\mu f(x)=\eta^{-1}(f(x+\eta e_\mu)-f(x)),\qquad
 (d\lambda)_\mu=\partial_\mu\lambda,\qquad
 (d\mathscr A)_{\mu\nu}
   =\partial_\mu\mathscr A_\nu-\partial_\nu\mathscr A_\mu.
\]
Write \(\delta=d^T\) and \(\Delta_{\mathrm{sc}}=\delta d\)
on scalars. Distances between sites are periodic sup distances.
For bonds use their initial sites, and for coordinate supports use
the minimum distance between their points.
The local regularity norms \(\|B\|_{i,W}\), \(i=1,2,3\), mean
the supremum, first-difference supremum, and \(\alpha\)-H\"older
first-difference supremum on \(W\), respectively. Only contained
bonds and difference cells enter these local norms.
Elsewhere \(\|B\|_2\) denotes the quadratic norm; the domain
section explicitly uses \(\|B\|_i\) for the three global
regularity norms.
H\"older differences compare parallel output bonds at distance
in \((0,1]\), with the input column and derivative direction fixed.

For \(z\in\Z^3/\ell\Z^3\), put
\(\Delta_z=\{c:|c-z|_\infty<1/2\}\).
The partition \(M\)-cubes have side \(M\) and centers in
\(M\Z^3/\ell\Z^3\).
The scalar, bond, and plaquette averages to the unit lattice are
\begin{align}
 (\mathbf Q_k\lambda)(z)
   &=\eta^3\sum_{c\in\Delta_z}\lambda(c),\notag\\
 (\mathbf Q_k\mathscr A)_\mu(z)
   &=\eta^3\sum_{c\in\Delta_z}\sum_{j=0}^{N-1}
          \eta\mathscr A_\mu(c+j\eta e_\mu),\label{eq:photon-block-averages}\\
 (\mathbf Q_k^{(2)}F)_{\mu\nu}(z)
   &=\eta^3\sum_{c\in\Delta_z}\sum_{i,j=0}^{N-1}
          \eta^2F_{\mu\nu}(c+i\eta e_\mu+j\eta e_\nu).\notag
\end{align}
All sums are taken on the periodic covering lattice.
Let \(\mathcal R_k\) be the orthogonal projection onto
\(\Delta_{\mathrm{sc}}(\ker\mathbf Q_k)\) in scalar cochains.
The Landau representative \(H_k^{\mathrm{ph}}A\) minimizes
\(\|d\mathscr A\|_2^2\) under
\begin{equation}\label{eq:photon-Landau-constraints}
 \mathbf Q_k\mathscr A=A,\qquad \mathcal R_k\delta\mathscr A=0.
\end{equation}
Set
\begin{equation}\label{eq:photon-coarse-form}
 \Delta_k=(dH_k^{\mathrm{ph}})^TdH_k^{\mathrm{ph}}.
\end{equation}

For one block step on the unit lattice, let
\(B(y)=\{z:|z-y|_\infty<L/2\}\), \(y\in L\Z^3/\ell\Z^3\).
The scalar average is
\(\mathbf Q\lambda(y)=L^{-3}\sum_{z\in B(y)}\lambda(z)\).
On bonds put
\[
 (\mathbf QA)(y,y+Le_\mu)
   =L^{-4}\sum_{z\in B(y)}A(\Gamma_{z,z+Le_\mu}),\qquad
 (\tau A)(y,z)=\tfrac16\sum_\pi A(\Gamma^\pi(y,z)).
\]
Here \(A(\Gamma)\) is the oriented bond sum; the first paths are
straight, while the latter change coordinates in each of the six
orders. Write \(\mathcal Z=\{Z:\mathbf QZ=0,\ \tau Z=0\}\).
For a coarse field \(A'\) on side \(\ell/L\), set
\(A'_L(b)=L^{-1/2}A'(b/L)\).

\begin{theorem}\label{thm:maxwell-block}
Assume that \(\ell/(LM)\) is an odd integer at least \(13\).
There are constants \(C_P^{\mathrm{ph}}>0\) and
\(0<\gamma_P\le1\), depending only on \(L,\alpha\), for which
the following assertions hold uniformly in \(k,\ell\).

The minimizer \(H_k^{\mathrm{ph}}\) exists uniquely, is real
linear, and satisfies \(\mathbf Q_kH_k^{\mathrm{ph}}=1\).
It preserves constants and closed fields, and \(H_0^{\mathrm{ph}}=1\).
Each of its kernel, first-difference kernel, and short H\"older
first-difference kernel is bounded by
\(C_P^{\mathrm{ph}}e^{-\gamma_Pd(b,b')}\).
Also
\[
 |\Delta_k(b,b')|\le C_P^{\mathrm{ph}}e^{-\gamma_Pd(b,b')}.
\]

There is a real coordinate space \(V_k\) with an orthonormal
basis \(\{\Upsilon\}\), each vector supported in an \(L\)-block
or on a noncentral face bond, and a bijection \(C:V_k\to\mathcal Z\).
The basis choices are translates of one fixed block basis, and
\[
 \dim V_k=(2-2L^{-3})\ell^3,\qquad
 \|z\|_2\le\|Cz\|_2\le\sqrt{5L^3}\|z\|_2.
\]
The form \(\widetilde\Delta_k=C^T\Delta_kC\) is strictly positive,
with upper and lower bounds depending only on \(L\), and
\[
 \big|\widetilde\Delta_k^{-1/2}(\Upsilon,\Upsilon')\big|
 \le C_P^{\mathrm{ph}}e^{-\gamma_Pd(\Upsilon,\Upsilon')}.
\]
Here the distance on the right is between basis supports.

For every real \(A'\), the minimum \(A_k^{\min}(A')\) of
\(\langle A,\Delta_kA\rangle\) under
\(\mathbf QA=A'_L,\ \tau A=0\) exists uniquely. With
\(\mathrm{Sq}_k=C\widetilde\Delta_k^{-1/2}\),
\begin{equation}\label{eq:main-block-split}
 \begin{split}
 &\langle A_k^{\min}(A')+\mathrm{Sq}_k\zeta,\,
   \Delta_k(A_k^{\min}(A')+\mathrm{Sq}_k\zeta)\rangle\\
 &\hspace{35mm}=\|dH_{k+1}^{\mathrm{ph}}A'\|_2^2+\|\zeta\|_2^2.
 \end{split}
\end{equation}
The operator at level \(k+1\) on the right is formed on side
\(\ell/L\) and mesh \(\eta/L\). The harmonic part of \(A'\)
is arbitrary. The probability law on the affine constraint space
with density proportional to
\(\exp(-\langle A,\Delta_kA\rangle/2)\), relative to its Euclidean
volume, is the image of the standard product Gaussian \(\mu_I\)
under \(A=A_k^{\min}(A')+\mathrm{Sq}_k\zeta\).
For every \(q\ge1\),
\[
 \mu_I\{\sup_\Upsilon|\zeta(\Upsilon)|\le q\}
   =e^{-\ell^3\delta\varepsilon^0(q)},\qquad
 0\le\delta\varepsilon^0(q)\le12e^{-q^2/2}.
\]

If \(\square\) is a partition \(M\)-cube and
\(\square^\dagger\) its concentric enlargement of side
\((2L+1)M\), then for every real \(A\) there is an independent
real gauge \(\nu\) on this nonwrapping enlargement such that
\[
 \|H_k^{\mathrm{ph}}A-d\nu\|_{i,\square^\dagger}
       \le C_B\|dA\|_\infty,\qquad i=1,2,3,
\]
where \(C_B\) depends only on \(L,M,\alpha\); the subscripts
\(i\) denote the three regularity norms.

Finally, let \(\ell_2=J\ell_1\), with \(J\) odd and both sides
as above, and use the same translated block basis. Periodic
pullback intertwines \(H_k^{\mathrm{ph}},\Delta_k,C,
\widetilde\Delta_k,(\widetilde\Delta_k+xI)^{-1}\) for \(x\ge0\),
\(\widetilde\Delta_k^{-1/2},\mathrm{Sq}_k\), and
\(A_k^{\min}\) on the two tori.
\end{theorem}

We develop the Abelian block argument of
\cite[Section~10]{QED3Manuscript}. The two gauge choices have
different uses. The Landau representative admits a Fourier
description with constants uniform in the fine mesh. The axial
slice has local coordinates on which the curvature energy is
uniformly coercive. The inverse square root of this coordinate
form gives the Gaussian variables in \eqref{eq:main-block-split}.

The uniformity of the Landau bounds requires cancellation at the
zero Fourier alias before the mesh tends to zero. After that
cancellation, the remaining alias sums converge on a fixed complex
strip. A differentiated summability estimate, using \(\alpha<1\),
permits contour shifting for all three kernels. For the
fluctuation form, exponential conjugation bounds the resolvent
uniformly; its integral gives locality of the positive inverse
square root. These arguments are proved in
\cref{subsec:photon-symbol-regularity,subsec:photon-fluctuation-root}.

The remaining sections provide the photon estimates used in a
nonlinear block integration. Curvature controls the reconstructed
field after independent local real gauge changes. Polynomial
weakening along parameter paths preserves exponential decay and
separates linked components; its exact bounds are stated in
\cref{prop:photon-weakened-bounds}. The field-domain theorem,
\cref{prop:photon-domain-preservation}, states the additional
charge inequalities and the margin for complex fluctuations.
Only the exact operators are asserted to commute with periodic
pullback. No fermion operator or interacting estimate is assumed.
